% Generated draft: requires mathematical and source review.
\documentclass[11pt]{article}
\usepackage[margin=25mm]{geometry}
\usepackage{fontspec}
\setmainfont{Latin Modern Roman}
\usepackage{amsmath,amssymb,mathtools,unicode-math}
\setmathfont{Latin Modern Math}
\usepackage{enumitem,needspace,xurl,hyperref}
\setlength{\parindent}{0pt}
\setlength{\parskip}{5pt}
\setlength{\emergencystretch}{3em}
\newcommand{\N}{\mathbb{N}}
\newcommand{\Z}{\mathbb{Z}}
\newcommand{\Q}{\mathbb{Q}}
\newcommand{\R}{\mathbb{R}}
\newcommand{\C}{\mathbb{C}}
\newcommand{\F}{\mathbb{F}}
\newcommand{\A}{\mathbb{A}}
\newcommand{\PP}{\mathbb{P}}
\newcommand{\E}{\mathbb{E}}
\newcommand{\eps}{\varepsilon}
\newcommand{\dd}{\,\mathrm{d}}
\begin{document}
\section*{UnsolvedMath ID: 700654}
\textbf{Source problem code:} catalogue entry
\subsection{Definitions and mathematical statement}
Under the customary convention that a commutative ring has an identity, let \(R\) be such a ring. For an integer \(n\geq 2\), write
\[
\operatorname{SL}_n(R)=\{A\in M_n(R):\det(A)=1\},
\qquad
E_n(R)=\left\langle I_n+r e_{pq}:r\in R,\ 1\leq p,q\leq n,\ p\neq q\right\rangle,
\]
where \(I_n\) is the identity matrix and \(e_{pq}\) is the matrix unit. The membership problem for this inclusion asks for an algorithm which, given an encoded matrix \(A\in\operatorname{SL}_n(R)\), terminates and decides whether \(A\in E_n(R)\).

For the usual nonnegative indexing of Quillen algebraic \(K\)-theory, write \(K_i(R)=\pi_i(\mathbf K(R))\) for \(i\geq 0\), where \(\mathbf K(R)\) is the connective Quillen algebraic \(K\)-theory spectrum of \(R\). Relative to an effective specification of a group by generators and relations, its word problem asks for an algorithm which terminates on every finite word in the specified generators and their inverses and decides whether that word represents the identity.

The supplied questions of G. A. Noskov are:
\begin{enumerate}
\item[(a)] Is the membership problem for \(E_n(R)\leq\operatorname{SL}_n(R)\) soluble when \(R\) is a commutative ring?
\item[(b)] Is the word problem for \(K_i(R)\) soluble when \(K_i\) denotes a Quillen \(K\)-functor and \(R\) is a commutative ring?
\end{enumerate}
The supplied statement does not specify an effective representation or computability hypotheses for \(R\), an effective presentation of \(K_i(R)\), the intended ranges of \(n\) and \(i\), or whether the algorithms must be uniform in the ring and the indices. It also does not explicitly require rings to have an identity. The conventions above explain the usual objects but do not settle these missing qualifications. Consequently, the source questions cannot yet be expressed as fully specified algorithmic assertions without further source clarification.
\subsection{Short English statement}
Noskov asks whether one can decide membership in the elementary matrix subgroup of a special linear group over a commutative ring, and whether one can decide when a word represents the identity in a Quillen algebraic \(K\)-group of such a ring. The source leaves the computational input, ring hypotheses, index ranges, and uniformity requirements unspecified.
\subsection{Sources}
\begin{itemize}
\item Local UnsolvedMath registry, permanent ID \(700654\), exact source problem code \texttt{catalogue entry}, title \emph{Kourovka Notebook Problem 10.39}. The supplied record labels the problem \emph{open}; this is a supplied registry classification, not an independently verified mathematical or literature assessment. The record was updated on \texttt{2026-10-16T14:00:00Z}.
\item Supplied source attribution: \emph{The Kourovka Notebook (Unsolved Problems in Group Theory)}, problem 10.39, proposed by G. A. Noskov in the 10th issue (1986). The background identifies the 21st issue edition (2026), \texttt{arXiv:1401.0300v48}. The supplied collection description names E. I. Khukhro and V. D. Mazurov as editors.
\item The cited Notebook text has not been independently consulted for this draft. No page number, explicit computational conventions, or separate literature-verification date is supplied.
\end{itemize}
\end{document}
